\section{模型假设}

为使模型可识别并与题目数据条件一致，作如下假设：
\begin{enumerate}
  \item 题目给出的玻璃类型标签和表面风化标签总体可靠，可作为监督学习和分组统计的基准。
  \item 当14种成分比例之和位于\([85\%,105\%]\)时，检测误差处于可接受范围；超出区间的样本不参与主要模型拟合。
  \item 空白项表示对应成分低于检测限或未检出。进行一般统计时按题意记为0；进行对数比变换时采用\todo{说明零值替代方法及替代量}。
  \item 同一玻璃类型在相近风化状态下具有可比较的总体成分变化规律；在缺少充分纵向观测时，可利用横截面统计规律估计风化前成分。
  \item 未知样品与已分类样品来源于相同或相近的总体分布，且采用一致的检测口径。
  \item 未被题目记录的埋藏环境、年代等因素视为随机扰动，其总体影响由误差项吸收。
\end{enumerate}
